%% kapa: exam=1 round=28 year=2017 period=1 form=A subject=economics subject_ko=경제학원론 seq=17 number=57
%% source: Q-Net 감정평가사 기출문제 (cst00309) artlSeq=5199603 file=extracted/1008/5199603/2017년 제28회 감정평가사 1차시험 문제지/1교시 A형.pdf
%% conversion: kordoc 4.24.51 -> tools/gen1.py (auto); review=proofread-sonnet; flags=none
\begin{question}{57}
X재에 부과되던 물품세가 단위당 $t$에서 $2t$로 증가하였다. X재에 대한 수요곡선은 우하향하는 직선이며, 공급곡선은 수평일 때 설명으로 옳은 것은?
\begin{choices}
  \item 조세수입이 2배 증가한다.
  \item 조세수입이 2배보다 더 증가한다.
  \item 자중손실(deadweight loss)의 크기가 2배 증가한다.
  \item 자중손실의 크기가 2배보다 더 증가한다.
  \item 새로운 균형에서 수요의 가격탄력성은 작아진다.
\end{choices}
\end{question}
\answer{4}
